\section{Appendix: Agent Instruction Context}
\label{app:instruction-context}

This appendix expands the instruction-stream contract in implementation-neutral
terms. The method does not require a particular repository layout. It requires
that each agent receive a bounded, ordered, versioned text context and that the
rendered context be hashed and stored with the resulting artifact.

\subsection{Shared Cross-Agent Protocol}

All agents begin with the same legal reasoning protocol. The protocol has five
rules.

First, source authority controls legal substance. The source case, statute,
Restatement section, regulation, contract, or other supplied authority is the
controlling legal input for the run. Doctrine examples are drafting aids only.

Second, calibration domains are not global assumptions. The reported validation
uses Statute of Frauds because it has useful boundary facts, but the agents must
infer doctrine from the source. If the source supports another doctrine, the same
packet, rubric, clustering, and judging architecture is used with
doctrine-appropriate gates, elements, defenses, and exceptions.

Third, the schema is doctrine-general. Fields such as
\texttt{doctrine\_gates}, \texttt{primary\_gate\_id},
\texttt{rule\_trigger}, and \texttt{exception\_or\_defense} refer to the
controlling decision points of the source-supported doctrine. They are not
Statute-of-Frauds-only labels.

Fourth, variations must identify the legal effect of the edited fact. Invariant
variations change legally irrelevant surface facts. Material variations change
facts that should affect a doctrine element, defense, exception, burden, remedy,
or outcome. Each variation records the expected behavior before model answers
are generated.

Fifth, uncertainty is preserved. If the source is too thin,
jurisdiction-specific, procedurally tangled, or unsupported by the current
instruction context, the agent must state the limitation and preserve it for
escalation rather than inventing missing law.

\subsection{Frank Context}

Frank's instruction context defines source-to-benchmark generation. It tells
Frank to perform source intake, infer the doctrine family, identify material
facts and source limits, generate a neutral benchmark question, produce a
source-grounded gold answer, create boundary variations, and freeze a controller
card for downstream stages. Frank should stop rather than force a benchmark when
the source cannot support a stable packet.

The source extraction component records authority level, jurisdiction,
procedural posture, legal issue, rule statement, trigger facts, holding or
best-supported answer path, limits, source gaps, benchmark confidence, and
review flags. The gold-packet component records doctrine family, controlling
trigger, gate order, supporting facts, non-trigger facts, competing barriers,
exceptions or substitutes, jurisdiction-sensitive points, likely model mistakes,
and benchmark posture. The controller card freezes the current question, gold
answer, primary gate id, trigger test, trigger facts, non-triggered sibling
gates, compliance status, strongest counterargument, fallback theories, omitted
control facts, and source metadata.

Frank's question-writing component controls what response models see. The
question must be neutral, exam-like, and source-faithful. It should preserve the
holding-driving facts and one or two plausible distractors while avoiding answer
leakage. The call of the question should not reveal the controlling doctrine,
tell the model which side should win, or turn the prompt into a checklist of
subissues. Jurisdiction is stated only when the benchmark depends on
jurisdiction-specific law.

\subsection{Karthic Context}

Karthic's instruction context defines rubric construction from the locked Frank
packet. Karthic is not asked to create a generic grading sheet. It builds a
modular scoring instrument derived from the gold answer, source-supported
doctrine, expected failure modes, and, when available, observed response
clusters.

Karthic first audits the packet fields and marks them as fixed, fixed but
jurisdiction-sensitive, or needing confirmation. It then locks the doctrinal path
by stating the controlling doctrine, gate order, benchmark posture, and strongest
counterargument. Rubric rows are organized into modules for metadata, structural
gatekeeping, primary doctrine gates, fallback doctrines and defenses, and
cross-cutting answer discipline. Scored rows require observable legal behavior,
not stylistic preferences.

Each full row card should contain a row code, module, weight, description,
non-applicability guidance, golden target summary, observable features of a
strong answer, allowed omissions, contradiction flags, comparison guidance,
scoring anchors, failure labels, and row status. Broad rows are split only when
the split improves legal discrimination. Overlapping rows require a distinctness
rationale so that one legal error is not counted twice. Weights are assigned by
legal importance rather than prose length, and answer-level penalties or caps
are applied only after row scoring.

\subsection{Dasha Context}

Dasha's instruction context defines reasoning-signature extraction and
centroid-level evaluation. Dasha does not write the question, gold answer, or
rubric. It reads raw model responses after generation and extracts a normalized
description of the legal path each response appears to use.

The Dasha track audit records which question version is active, which rubric
belongs to that question, whether the run is single-track or perturbation-track,
which judge model is used, whether judge aggregation is single-model or
ensemble-based, and whether the judge is separate from the response-model pool.
Missing inputs that would distort scoring are review flags rather than fields to
silently invent.

Dasha's primary scored unit is the centroid. A centroid receives row-level
scores under the Karthic rubric. If perturbation tracks are active,
original-question outputs are clustered and scored only against the base track,
and variation outputs are clustered and scored only against their corresponding
variation track. Comparisons between original and variation tracks occur only
after both tracks have been independently clustered.

Dasha also records centroid composition metadata: cluster size, per-model answer
counts and shares, represented model count, dominant model, and dominant model
share. This metadata is descriptive and must not become a scoring input.

\subsection{Judge and Zak Context}

The judge context requires row-level rubric application. For every cluster
representative, the judge receives the representative response, the cluster
legal signal, and the Karthic rubric rows. It returns exactly one score and
rationale for every row. The score scale, weighting rule, and aggregation rule
are fixed in the run configuration.

Zak is the escalation layer. It creates packets only when the frozen thresholds
indicate uncertainty, disagreement, mixed clusters, missing rows, or other
configured failure states. An escalation packet contains the active track,
disputed clusters, relevant rubric rows, score summaries, response excerpts,
reason for escalation, and a concise handoff note. If no escalation is required,
the run still records an explicit no-escalation state.
